\documentclass[11pt,a4paper]{article}
\usepackage[a4paper,margin=0.85in]{geometry}
\usepackage{amsmath,amssymb,booktabs,array}
\usepackage[T1]{fontenc}
\usepackage{lmodern}
\usepackage{xcolor}
\usepackage{hyperref}
\usepackage{listings}
\usepackage{enumitem}
\hypersetup{colorlinks=true,linkcolor=blue,urlcolor=blue,citecolor=blue}
\lstset{basicstyle=\ttfamily\small,breaklines=true,frame=single,backgroundcolor=\color{gray!8}}

\title{Intermediate Task 1 Report\\\large Goal-Conditioned Treasure Hunt using Universal Value Function Approximators}
\author{AVD894 Reinforcement Learning Course Project}
\date{November 2026}

\begin{document}
\maketitle

\begin{abstract}
This report describes the first milestone of our course project based on
Universal Value Function Approximators (UVFAs) by Schaul et al. The paper's
central idea is to estimate the value of a state relative to a specified goal,
so that one value model can be reused for many goals. For Intermediate Task 1,
we formulate a small goal-conditioned treasure-hunt game as a finite Markov
Decision Process (MDP) and solve it exactly using value iteration. The exact
solution serves as a mathematical reference for the simulator and the learning
agent used in Intermediate Task 2.
\end{abstract}

\section{Paper in a nutshell}

The paper ``Universal Value Function Approximators'' introduces value functions
of the form
\[
    V(s,g;\theta),
\]
where $s$ is the current state and $g$ is a goal descriptor. A conventional
value function is tied to one objective, whereas a UVFA is intended to represent
many goal-conditioned value functions in one parameterized model. The model can
share useful structure across both states and goals and can therefore estimate
values for goals that were not explicitly observed during training.

The paper develops state and goal representations and combines their learned
embeddings. It also explains how the representation can be updated from
reinforcement-learning experience. The contribution is important because it
connects value-function approximation with multi-goal and transfer settings:
the same agent can be asked to reach different targets without training an
entirely independent value function for every target.

The primary source is the ICML 2015 paper by Schaul, Horgan, Gregor, and Silver,
available at \url{https://proceedings.mlr.press/v37/schaul15.html}.

\section{Our project idea}

We adapt the UVFA idea into a small, visual game called \emph{Goal-Conditioned
Treasure Hunt}. The player is placed in a maze and must reach a specified goal
cell. The goal may change between episodes. The project asks whether a single
goal-conditioned value model can make useful decisions for many target cells.

For Task 1, the purpose is not yet to train a neural network. Instead, we create
a small finite MDP and calculate its optimal value function exactly. This gives
us a ground-truth reference for later comparison. The adaptation is explicitly
labelled as a course-project environment; it is not claimed to be an exact copy
of every environment used in the paper.

\section{Mathematical problem and MDP formulation}

Let the maze contain a finite set of free cells $\mathcal{C}$. A state contains
the current player position $p$ and the selected goal position $g$:
\[
    s=(p,g), \qquad p,g\in\mathcal{C}, \quad p\neq g.
\]
The action space is
\[
    \mathcal{A}=\{\text{up},\text{down},\text{left},\text{right}\}.
\]

The transition is deterministic. A legal action moves the player to an adjacent
free cell. An action into a wall leaves the player in the same cell. The reward
used in Task 1 is
\[
r(s,a)=
\begin{cases}
+10, & \text{if the action reaches the goal},\\
-0.1, & \text{if the action is a legal non-terminal movement},\\
-1, & \text{if the action attempts to enter a wall}.
\end{cases}
\]
The episode terminates when $p=g$. The objective is to maximize the discounted
return
\[
    J(\pi)=\mathbb{E}_{\pi}\left[\sum_{t=0}^{T-1}\gamma^t r_t\right],
\]
with $\gamma=0.95$.

For a deterministic transition function $T(s,a)$, value iteration applies the
Bellman optimality update
\[
    V_{k+1}(s)=\max_{a\in\mathcal{A}}
    \left[r(s,a)+\gamma V_k(T(s,a))\right].
\]
The continuation value is set to zero when the next state is terminal. The
greedy policy is
\[
    \pi(s)=\arg\max_{a\in\mathcal{A}}
    \left[r(s,a)+\gamma V(T(s,a))\right].
\]

\section{Implementation of the paper idea}

The Task 1 implementation uses the project map in
\texttt{configs/task1\_small.json}. It contains 22 free cells. Because both the
player position and the goal position are part of the state, the number of
ordered state-goal pairs is
\[
    22\times21=462.
\]
The state space is therefore larger than the number of visible cells, while the
action space remains four-dimensional and discrete.

The implementation proceeds as follows:
\begin{enumerate}[leftmargin=*]
    \item Load the fixed maze and reward parameters.
    \item Enumerate every valid ordered pair of distinct free cells.
    \item Initialize all values to zero.
    \item Evaluate all four actions for every state.
    \item Apply the Bellman optimality update synchronously.
    \item Stop when the maximum value change is below the configured tolerance.
    \item Extract the greedy policy and follow it from a selected start to a goal.
\end{enumerate}

This is a small exact reference solution rather than a neural approximation. It
is useful because it confirms the MDP definition before we introduce replay,
target networks, and function approximation in Task 2.

\section{Code in a nutshell}

The relevant source files are:
\begin{itemize}[leftmargin=*]
    \item \texttt{src/goal\_treasure/env.py}: maze dynamics, actions, rewards, and observations.
    \item \texttt{src/goal\_treasure/exact.py}: value iteration and policy rollout.
    \item \texttt{scripts/run\_task1.py}: command-line entry point that writes JSON and CSV outputs.
    \item \texttt{configs/task1\_small.json}: reproducible map and numerical settings.
\end{itemize}

The main command is:
\begin{lstlisting}[language=bash]
python scripts/run_task1.py --config configs/task1_small.json \\
    --output artifacts/task1
\end{lstlisting}

The command writes a summary file and a policy table. The same calculation can
also be run from \texttt{notebooks/AVD894\_Task1\_Task2.ipynb}.

\section{Results}

The measured Task 1 result was obtained on 2 November 2026 using the supplied
configuration:

\begin{table}[h]
\centering
\caption{Exact Task 1 value-iteration result.}
\begin{tabular}{lr}
\toprule
Quantity & Measured value\\
\midrule
Free cells & 22\\
Ordered state-goal pairs & 462\\
Number of actions & 4\\
Discount factor & 0.95\\
Value-iteration iterations & 10\\
Final maximum value change & 0.0\\
Example path length & 9 moves\\
\bottomrule
\end{tabular}
\end{table}

For the example start at $(1,1)$ and goal at $(5,6)$, the exact policy produced
the path
\[
(1,1)\rightarrow(2,1)\rightarrow(3,1)\rightarrow(4,1)
\rightarrow(5,1)\rightarrow(5,2)\rightarrow(5,3)
\rightarrow(5,4)\rightarrow(5,5)\rightarrow(5,6).
\]

The result shows that the simplified problem is correctly specified and that a
policy can be extracted from the value function. It does not yet measure neural
generalization; that is the focus of Task 2.

\section{Limitations and next step}

The Task 1 map is fixed and small, and the transition is deterministic. The
solution is exact because every state-goal pair can be enumerated. This does not
scale directly to large maps or continuous observations. In Task 2, we therefore
replace the table with a neural goal-conditioned value model and evaluate it on
randomly generated start/goal episodes.

\section{References}

\begin{thebibliography}{9}
\bibitem{schaul2015}
Tom Schaul, Daniel Horgan, Karol Gregor, and David Silver,
``Universal Value Function Approximators,''
\emph{Proceedings of the 32nd International Conference on Machine Learning},
PMLR 37, pp. 1312--1320, 2015.
\url{https://proceedings.mlr.press/v37/schaul15.html}
\end{thebibliography}

\end{document}

